\documentclass[a4paper,11pt]{article}
\usepackage[utf8]{inputenc}
\usepackage[T1]{fontenc}
\usepackage[french]{babel}
\usepackage{lmodern}
\usepackage[margin=2.2cm]{geometry}
\usepackage{amsmath,amssymb}
\usepackage{booktabs}
\usepackage{enumitem}
\usepackage{xcolor}
\usepackage{tikz}

\setlength{\parindent}{1.5em}
\setlength{\parskip}{0.3em}

\title{Mod\'elisation CSP de l'\'enum\'eration des Kekul\'e et de la maximisation des couvertures de Clar}
\author{Toky Nirina RAMAHERISOA \\ \small{Stage M2 IMD \textendash{} 2025/2026}}
\date{\today}

\begin{document}
\maketitle

\section{Contexte et port\'ee}

Ce document pr\'esente un \textbf{mod\`ele CSP} (Programmation Par Contraintes) unifi\'e
pour deux probl\`emes centraux de l'analyse topologique des hydrocarbures aromatiques polycycliques
(HAP) : (i) l'\'enum\'eration des \emph{structures de Kekul\'e} d'une mol\'ecule, et (ii) la
\emph{maximisation du nombre de sextets de Clar} (couvertures de Clar). Les deux probl\`emes se
formulent naturellement dans un m\^eme cadre, ce qui permet une m\'ethode unique et d\'eclarative,
facilement extensible.

\paragraph{Position par rapport \`a l'impl\'ementation.}
L'impl\'ementation \emph{actuelle} dans le module \texttt{molviz} repose sur des algorithmes
combinatoires d\'edi\'es (Edmonds blossom pour le matching maximum, backtracking canonique pour
\'enum\'erer les Kekul\'e, \'enum\'eration exhaustive sur les sous-ensembles d'hexagones pour
Clar). Ces approches sont efficaces sur la fen\^etre $h_3$\textendash{}$h_9$ qui nous int\'eresse
et \'evitent toute d\'ependance \`a un solveur externe. Le mod\`ele CSP d\'ecrit ici a une
finalit\'e \textbf{m\'ethodologique} :
\begin{itemize}
  \item exposer la structure d\'eclarative commune aux deux probl\`emes ;
  \item s'aligner avec les outils existants (\texttt{pah-tools}, BenzAI~\cite{varet2023benzai})
    qui adoptent une mod\'elisation PPC pour Kekul\'e ;
  \item ouvrir la voie \`a l'ajout de contraintes additionnelles (sym\'etries, contraintes
    chimiques de stabilit\'e, etc.) sans toucher au noyau combinatoire.
\end{itemize}

\section{Notations et pr\'eliminaires}

Soit $G = (V, E)$ le squelette carbone d'une mol\'ecule HAP : $V$ est l'ensemble des atomes de
carbone (les hydrog\`enes sont retir\'es par construction), $E$ l'ensemble des liaisons $C\!-\!C$.
On note $|V| = n$, $|E| = m$. Pour $v \in V$, $\delta(v) \subseteq E$ d\'esigne l'ensemble des
ar\^etes incidentes \`a $v$.

Soit $\mathcal{H} = \{h_1, \dots, h_k\}$ l'ensemble des \emph{cycles hexagonaux} de $G$. Pour
$h \in \mathcal{H}$, on note $V(h) \subset V$ ses 6 atomes et $E(h) \subset E$ ses 6 ar\^etes.

\begin{definition}[Structure de Kekul\'e]
Une \emph{structure de Kekul\'e} de $G$ est un couplage parfait $M \subseteq E$, c'est-\`a-dire
un sous-ensemble d'ar\^etes (les ``doubles liaisons'') tel que chaque sommet est dans exactement
une ar\^ete de $M$. Pour une mol\'ecule sans couplage parfait (typiquement \`a nombre impair de
carbones ou topologie radicalaire), on consid\`ere \`a la place un \emph{matching maximum},
laissant un certain nombre de sommets non couverts qui sont les \emph{sites radicalaires}.
\end{definition}

\begin{definition}[Couverture de Clar]
Une \emph{couverture de Clar} (Clar 1972) est un choix d'un sous-ensemble $S \subseteq \mathcal{H}$
d'hexagones \emph{deux-\`a-deux vertex-disjoints} (i.e.\ $V(h) \cap V(h') = \emptyset$ pour tout
$h \neq h' \in S$), accompagn\'e d'un matching parfait du sous-graphe induit par les atomes non
consomm\'es par $S$. Chaque hexagone de $S$ porte un \emph{sextet aromatique} (rond de Clar)
de 6 \'electrons~$\pi$ d\'elocalis\'es. Le \emph{nombre de Clar} de la mol\'ecule est
$\max |S|$ sur toutes les couvertures valides.
\end{definition}

Seuls les hexagones peuvent porter un sextet (r\`egle de H\"uckel $4n+2$ avec $n=1$~: il faut
exactement 6 \'electrons~$\pi$). Les pentagones et heptagones sont donc structurellement exclus
de $\mathcal{H}$ dans notre mod\'elisation.

\section{Mod\`ele CSP $\mathcal{M}_K$ pour les structures de Kekul\'e}

\paragraph{Variables.} Une variable bool\'eenne par ar\^ete~:
\[
  x_e \in \{0, 1\}, \quad \forall e \in E,
\]
avec l'interpr\'etation $x_e = 1 \iff e$ est une double liaison.

\paragraph{Contraintes.} Chaque atome porte exactement une double liaison~:
\begin{equation}
  \sum_{e \in \delta(v)} x_e = 1, \quad \forall v \in V.
  \tag{C-Kek}
\end{equation}

\paragraph{Objectif.} L'\'enum\'eration des solutions de $\mathcal{M}_K$ produit l'int\'egralit\'e
des structures de Kekul\'e de la mol\'ecule. Le nombre de solutions est exactement le
\emph{nombre de Kekul\'e} (qui pour les benz\'eno\"ides peut \^etre calcul\'e en temps
polynomial par le th\'eor\`eme de Rispoli 2001, cas particulier du th\'eor\`eme de Kasteleyn 1967
pour les graphes planaires).

\paragraph{Extension radicalaire.} Si $G$ n'admet aucun couplage parfait, on relaxe l'\'egalit\'e
en in\'egalit\'e et on maximise la cardinalit\'e~:
\begin{equation}
  \sum_{e \in \delta(v)} x_e \leq 1 \quad \forall v \in V,
  \qquad \text{maximiser} \quad \sum_{e \in E} x_e.
\end{equation}
Les sommets $v$ tels que $\sum_{e \in \delta(v)} x_e = 0$ \`a l'optimum sont les sites
radicalaires.

Ce mod\`ele est exactement celui utilis\'e par \texttt{pah-tools}~\cite{pahtools} (cf.\
fichier \texttt{kekule.py}, qui d\'efinit le m\^eme syst\`eme via PyCSP3 puis appelle
\texttt{solve(sols=ALL)}). Il correspond au mod\`ele $\mathcal{M}_{1K}$ de la th\`ese
de Varet~\cite{varet2023benzai} (section~6.3.1, comparaison \'etablie avec la m\'ethode
de Rispoli).

\section{Mod\`ele CSP $\mathcal{M}_C$ pour les couvertures de Clar}

L'\'el\'egance de la formulation CSP appara\^it ici~: on \'etend simplement $\mathcal{M}_K$ avec
un jeu de variables suppl\'ementaires pour les sextets, et une \emph{unique contrainte de
couplage} qui int\`egre simultan\'ement la vertex-disjonction et l'existence du matching
parfait du r\'esidu.

\paragraph{Variables.} Aux variables d'ar\^etes $x_e$, on ajoute une variable bool\'eenne par
hexagone~:
\[
  y_h \in \{0, 1\}, \quad \forall h \in \mathcal{H},
\]
avec l'interpr\'etation $y_h = 1 \iff h$ porte un sextet de Clar (rond de Clar).

\paragraph{Contrainte unifi\'ee.} Pour chaque atome $v$, exactement \emph{une} des deux
situations suivantes se produit~: soit $v$ appartient \`a un sextet (et alors aucune de ses
ar\^etes n'est double, car ses 4 \'electrons sont engag\'es dans le sextet aromatique), soit
$v$ est dans le r\'esidu et doit \^etre couvert par une double liaison~:
\begin{equation}
  \sum_{e \in \delta(v)} x_e \;+\; \sum_{h \in \mathcal{H} \,:\, v \in V(h)} y_h \;=\; 1,
  \quad \forall v \in V.
  \tag{C-Clar}
\end{equation}

\paragraph{Objectif.} On maximise le nombre de sextets~:
\begin{equation}
  \text{maximiser} \quad Z = \sum_{h \in \mathcal{H}} y_h.
\end{equation}

\paragraph{Pourquoi cette contrainte capte tout.} Examinons les cas~:
\begin{itemize}
  \item \textbf{Vertex-disjonction des sextets.} Si deux hexagones $h, h' \in S$ partageaient un
    atome $v \in V(h) \cap V(h')$, on aurait $y_h + y_{h'} \geq 2$ et donc le membre de gauche
    de (C-Clar) en $v$ serait $\geq 2$, violant l'\'egalit\'e \`a 1.
  \item \textbf{Atomes consomm\'es par un sextet.} Si $v$ appartient \`a un unique sextet $h$
    (ce qui est forc\'e par le point pr\'ec\'edent), alors $\sum y_h = 1$ et la contrainte
    impose $\sum x_e = 0$~: aucune ar\^ete incidente \`a $v$ n'est double dans le mod\`ele.
    Les 4 \'electrons de valence de $v$ se r\'epartissent : 3 dans les $\sigma$-liaisons aux
    voisins, 1 dans le syst\`eme $\pi$ aromatique du sextet.
  \item \textbf{Matching parfait du r\'esidu.} Si $v$ n'appartient \`a aucun sextet, alors
    $\sum y_h = 0$ et la contrainte impose $\sum x_e = 1$~: $v$ est couvert par exactement une
    double liaison. C'est la d\'efinition d'un matching parfait sur les atomes du r\'esidu.
\end{itemize}

\paragraph{Cas particulier.} Si on impose $y_h = 0$ pour tout $h \in \mathcal{H}$, le mod\`ele
$\mathcal{M}_C$ se r\'eduit \`a $\mathcal{M}_K$. Les solutions sont alors les structures de
Kekul\'e classiques. Ainsi, $\mathcal{M}_K$ est un sous-mod\`ele strict de $\mathcal{M}_C$.

\paragraph{Extension radicalaire.} Si $G$ ne poss\`ede pas suffisamment de structure pour que
chaque atome soit dans un sextet ou dans une double, on relaxe (C-Clar) en in\'egalit\'e et on
introduit un objectif lexicographique~:
\begin{equation}
  \sum_{e \in \delta(v)} x_e + \sum_{h \,:\, v \in V(h)} y_h \;\leq\; 1, \quad \forall v \in V,
\end{equation}
\begin{equation}
  \text{maximiser lex.} \quad
  \Bigl( -\!\!\sum_{v \in V}\! r_v, \;\; \sum_{h} y_h \Bigr),
  \quad \text{o\`u} \quad r_v = 1 - \!\!\sum_{e \in \delta(v)}\! x_e - \!\!\sum_{h \,:\, v \in V(h)}\! y_h
\end{equation}
$r_v$ vaut 1 si $v$ est un radical, 0 sinon. La maximisation lexicographique \emph{minimise
d'abord le nombre total de radicaux} (qui vaut alors la d\'eficience globale du graphe),
\emph{puis maximise le nombre de sextets} sous cette contrainte. Cela rend la d\'efinition
op\'erationnelle pour les mol\'ecules \`a parit\'e impaire ou topologie radicalaire (typiques
des structures 5/7 \'etudi\'ees dans notre stage).

\section{Complexit\'e et r\'esolution pratique}

\subsection*{Taille du mod\`ele $\mathcal{M}_C$}

\begin{center}
\begin{tabular}{lcc}
  \toprule
  Composante & Nombre & Pour $h_9$ typique \\
  \midrule
  Variables $x_e$ & $|E|$ & $\sim 40$\textendash{}$50$ \\
  Variables $y_h$ & $|\mathcal{H}|$ & $\leq 9$ \\
  Contraintes (C-Clar) & $|V|$ & $\sim 30$\textendash{}$35$ \\
  \bottomrule
\end{tabular}
\end{center}

\subsection*{Choix algorithmiques}

Le mod\`ele $\mathcal{M}_C$ se r\'esout efficacement avec des solveurs CP standard (Choco, ACE
via PyCSP3). L'\'enum\'eration de toutes les couvertures \emph{optimales} (i.e.\ avec
$Z = \max$) demande un appel \`a la recherche \`a deux passes : (1) calcul de $Z^\star$ via
\texttt{minimize}/\texttt{maximize}, (2) ajout de la contrainte $\sum y_h = Z^\star$ puis
\texttt{searchAllSolutions}. La premi\`ere passe \'etablit le nombre de Clar~; la seconde
\'enum\`ere les couvertures qui l'atteignent.

\subsection*{Lien avec notre impl\'ementation \texttt{molviz}}

Notre impl\'ementation actuelle r\'esout le m\^eme probl\`eme par \'enum\'eration directe~:

\begin{center}
\begin{tabular}{lll}
  \toprule
  & Mod\`ele $\mathcal{M}_C$ (CSP) & Impl\'ementation \texttt{clar.py} \\
  \midrule
  Variables & $x_e$ + $y_h$ & Bitmask sur $\mathcal{H}$ \\
  R\'esolution & Solveur PPC (Choco/ACE) & Backtracking + Edmonds blossom \\
  Compl\'exit\'e & D\'epend du solveur & $O(2^{|\mathcal{H}|} \cdot n^3)$ \\
  D\'ependance & PyCSP3 + JVM & Pur Python (NetworkX) \\
  Pour $|\mathcal{H}| \leq 9$ & $<1$~s & $<10$~ms \\
  \bottomrule
\end{tabular}
\end{center}

\noindent Les deux approches sont \'equivalentes en termes de r\'esultats (m\^eme ensemble de
couvertures optimales). L'\'enum\'eration directe est marginalement plus rapide sur nos petits
graphes en raison du faible $|\mathcal{H}|$ et de l'absence d'overhead JVM. Le mod\`ele CSP
devient pr\'ef\'erable d\`es que l'on souhaite \emph{combiner} d'autres contraintes
(sym\'etries, motifs interdits, contraintes de stabilit\'e thermique, etc.), ce qui est
naturellement extensible dans le formalisme PPC mais demanderait un refactor cons\'equent dans
l'impl\'ementation actuelle.

\section{Discussion}

Le mod\`ele $\mathcal{M}_C$ illustre la puissance d\'eclarative de la PPC pour les probl\`emes
de chimie topologique~: une \emph{unique} contrainte par atome capte simultan\'ement la
vertex-disjonction des sextets, le matching parfait du r\'esidu, et la s\'eparation entre les
\'electrons engag\'es dans les $\sigma$ et ceux d\'elocalis\'es dans les $\pi$. Cette concision
contraste avec l'\'enum\'eration combinatoire \`a deux niveaux (sous-ensembles d'hexagones~+
matching du r\'esidu) qu'on \'ecrit dans une approche algorithmique directe.

Pour les besoins de notre stage (analyse des structures $h_3$\textendash{}$h_9$ avec
\'eventuellement des cycles 5 et 7), l'impl\'ementation directe est suffisante et plus l\'eg\`ere
en d\'ependances. Le mod\`ele CSP reste pertinent comme \textbf{r\'ef\'erence m\'ethodologique}
et comme \textbf{point d'entr\'ee} pour des \'evolutions ult\'erieures (combinaison avec les
contraintes du g\'en\'erateur de mol\'ecules, qui est d\'ej\`a en PPC, et qui pourrait b\'en\'eficier
de l'agr\'egation des contraintes topologiques et aromatiques dans un m\^eme solveur).

\begin{thebibliography}{9}
\bibitem{varet2023benzai}
Yannis Varet. \emph{G\'en\'eration exhaustive et caract\'erisation des hydrocarbures aromatiques
polycycliques benz\'eno\"ides par programmation par contraintes}. Th\`ese de doctorat,
Aix-Marseille Universit\'e, 2023.
\bibitem{pahtools}
\texttt{pah-tools} \textendash{} A collection of tools for processing Polycyclic Aromatic
Hydrocarbons. \url{https://github.com/<lab>/pah-tools}.
\end{thebibliography}

\end{document}
